% ============================================================================
% CopilotScope -- podręcznik: czym jest, jak mierzy i jak go używać.
%
% To NIE jest artykuł naukowy. Formalnym opracowaniem modelu oceny jest
% research/articles/quality_measurement_framework.tex; ten dokument to
% podręcznik, który czyta się po to, żeby zrozumieć system, a potem go
% obsługiwać -- i właśnie dlatego zawiera samouczek.
%
% Źródła prawdy, które ten plik prezentuje, a nie rozszerza: README.md,
% docs/TUTORIAL.md, docs/SIGNAL_COVERAGE.md, docs/PRIVACY.md,
% docs/CALIBRATION.md, SECURITY.md i GOVERNANCE.md. Nie wprowadza żadnego
% faktu, którego tamte dokumenty już nie zawierają; jeśli się z nimi rozejdzie,
% to one mają rację, a ten plik jest nieaktualny.
%
% Diagramy nie są tu rysowane po raz drugi. Powstają z docs/diagrams/*.mmd --
% tych samych plików, które osadza docs/DIAGRAMS.md, utrzymywanych w zgodzie
% przez scripts/check-diagrams.mjs -- renderowane do wektorowych PDF-ów przez
% scripts/render-diagrams.mjs. Jedno źródło, dwa formaty wyjściowe, a oba
% wydania językowe współdzielą wynik.
%
% Wydanie angielskie leży obok:
%   copilotscope-manual.tex
% scripts/check-doc-parity.mjs przerywa build, jeśli jedno zostanie zmienione
% bez drugiego.
%
% Budowanie (PDF to wynik builda -- nigdy nie jest commitowany):
%   npm ci && npm run render:diagrams
%   pdflatex copilotscope-manual.pl.tex && pdflatex copilotscope-manual.pl.tex
% Budowane też przez .github/workflows/build-docs-pdf.yml, który buduje oba
% wydania.
% ============================================================================
\documentclass[11pt,a4paper]{article}

\usepackage[utf8]{inputenc}
\usepackage[T1]{fontenc}
\usepackage[polish]{babel}
\usepackage{lmodern}
\usepackage{geometry}
\usepackage{xcolor}
\usepackage{hyperref}
\usepackage{graphicx}
\usepackage{enumitem}
\usepackage{fancyhdr}
\usepackage{titlesec}
\usepackage{parskip}
\usepackage{booktabs}
\usepackage{array}
\usepackage{tabularx}
\usepackage{longtable}
\usepackage{fancyvrb}
\usepackage{microtype}
\usepackage{amssymb}
\usepackage{url}
\urlstyle{tt}

\geometry{a4paper, left=2.2cm, right=2.2cm, top=2.5cm, bottom=2.5cm}

% Kolory domowe
\definecolor{primaryblue}{RGB}{41,128,185}
\definecolor{darkblue}{RGB}{31,78,121}
\definecolor{lightgray}{RGB}{236,240,241}
\definecolor{darkgray}{RGB}{52,73,94}
\definecolor{accentgreen}{RGB}{39,174,96}
\definecolor{accentorange}{RGB}{230,126,34}
\definecolor{accentred}{RGB}{231,76,60}

\titleformat{\section}
    {\color{primaryblue}\Large\bfseries\sffamily}
    {\thesection.}{0.5em}{}[\titlerule]
\titleformat{\subsection}
    {\color{darkblue}\large\bfseries\sffamily}
    {\thesubsection}{0.5em}{}
\titleformat{\subsubsection}
    {\color{darkgray}\normalsize\bfseries\sffamily}
    {\thesubsubsection}{0.5em}{}

\newcommand{\docstatus}{BEZ KALIBRACJI --- PATRZ ROZDZIAŁ 7}

\pagestyle{fancy}
\fancyhf{}
\renewcommand{\headrulewidth}{0.4pt}
\fancyhead[L]{\small\sffamily\color{darkgray} CopilotScope --- podręcznik}
\fancyhead[R]{\small\sffamily\color{accentred} \docstatus}
\fancyfoot[C]{\thepage}

\hypersetup{
    colorlinks=true,
    linkcolor=primaryblue,
    urlcolor=primaryblue,
    pdfauthor={Konrad Cinkusz},
    pdftitle={CopilotScope: ocena jako\'{s}ci sesji z asystentami AI}
}

% Kod w tekście. \detokenize czyni każdy znak literalnym, więc podkreślenie w
% kluczu konfiguracyjnym nie wymaga ucieczki w miejscu użycia -- to najczęstszy
% sposób, w jaki taki dokument przestaje się budować.
\newcommand{\code}[1]{\texttt{\detokenize{#1}}}
% Ścieżka lub długi identyfikator w wąskiej kolumnie tabeli. \path łamie po
% / . _ oraz - , czego \texttt nie robi.
\newcommand{\codepath}[1]{{\small\path{#1}}}
% To samo w rozmiarze tekstu, dla długiego identyfikatora w zdaniu.
\newcommand{\codewrap}[1]{\path{#1}}
\newcommand{\finding}[1]{{\color{primaryblue}\textbf{#1}}}

\DefineVerbatimEnvironment{shell}{Verbatim}
    {frame=single, rulecolor=\color{darkgray}, framesep=5pt,
     fontsize=\small, xleftmargin=2pt, xrightmargin=2pt}

\begin{document}
% Polski babel czyni znak " aktywnym. W tym dokumencie cudzysłowy zapisujemy
% wprost jako znaki Unicode, a " pojawia się w cytowanych fragmentach JSON,
% więc skrót jest wyłączony.
\shorthandoff{"}

\thispagestyle{plain}
\begin{center}
    {\LARGE\bfseries\sffamily\color{darkblue} CopilotScope}

    \vspace{0.3cm}

    {\Large\sffamily\color{darkgray} Ocena jako\'{s}ci sesji z asystentami
    programistycznymi AI --- co mierzy, jak dzia\l{}a i jak go uruchomi\'{c}}

    \vspace{0.6cm}

    {\large Konrad Cinkusz}

    \vspace{0.2cm}

    {\normalsize\color{darkgray} \url{https://github.com/konradcinkusz/copilot-scope}}

    \vspace{0.2cm}

    {\normalsize\color{accentred}\sffamily \docstatus\ --- \today\ ---
     \textit{English edition: } \texttt{copilotscope-manual.tex}}
\end{center}

\vspace{0.4cm}

\begin{center}
\begin{minipage}{0.92\linewidth}
\small
\textbf{Czym jest ten dokument.} Wszystkim, czego potrzeba, żeby zrozumieć
CopilotScope, a potem go obsługiwać: problemem, który adresuje, tym jak jest
zbudowany, tym co dokładnie mierzy i czego mierzyć odmawia, oraz samouczkiem,
który prowadzi czytelnika od pustej maszyny do ocenionej sesji, wdrożenia
zespołowego i panelu w Grafanie.

\textbf{Czym nie jest.} Nie jest artykułem naukowym --- tym jest
\texttt{research/articles/quality\_measurement\_framework.tex} --- ani
walidacją oceny. Wskaźnik złożony nie został skalibrowany względem ocen
ludzkich. Rozdział~7 mówi o tym wprost, ponieważ liczba podana bez swoich
ograniczeń jest gorsza niż brak liczby.
\end{minipage}
\end{center}

\vspace{0.5cm}

{\small\tableofcontents}

\newpage

% ============================================================================
\section{Problem i to, co jest mierzone zamiast niego}
% ============================================================================

Każdy dostawca asystenta programistycznego AI powie Ci, ile Twój zespół go
zużył. Stanowiska, pokazane podpowiedzi, przyjęte podpowiedzi, spalone tokeny.
Te liczby są prawdziwe, łatwe do zebrania i odpowiadają na pytanie, którego
nikt faktycznie nie zadaje. Pytanie, które zespół ma naprawdę, brzmi: czy to
narzędzie pomaga --- a zużycie nie może na nie odpowiedzieć, bo rośnie
dokładnie tak samo wtedy, gdy asystent działa dobrze, i wtedy, gdy programista
utknął w pętli, pytając cztery razy o to samo innymi słowami.

\finding{CopilotScope ocenia sesje, nie zużycie.} Sesja to jedna rozmowa z
asystentem. Dla każdej produkuje złożony wskaźnik jakości wraz z miarą pewności
obok niego, konkretną turę, w której rozmowa się zacięła, i powód, informację
czy przyjęte przez programistę zmiany faktycznie przetrwały w pliku, oraz koszt
pętli naprawczej w tokenach i czasie.

Trzy własności kształtują wszystko, co następuje dalej:

\begin{enumerate}[leftmargin=*]
    \item \textbf{Nic nie opuszcza maszyny.} Kolektor działa na Twojej
    infrastrukturze i nie ma żadnego wywołania zwrotnego. To nie jest deklaracja
    polityki; w kodzie nie ma ścieżki, która gdziekolwiek dzwoni.
    \item \textbf{Wzór jest jawny i deterministyczny.} Sześć ważonych składowych,
    renormalizowanych po tych, które faktycznie dostarczyły dane, z miarą
    pewności eksportowaną obok każdej oceny. Każdą liczbę można wyprowadzić
    ręcznie.
    \item \textbf{Nie da się nim rankingować programistów.} W żadnym widoku ani
    eksporcie nie istnieje wymiar „per programista'', a testy tego pilnują. To
    jest egzekwowane, a nie obiecane --- rozdział~7.3 wyjaśnia, dlaczego to
    ograniczenie jest nośne, a nie ozdobne.
\end{enumerate}

\subsection{Dla kogo to jest}

Dla lidera zespołu, który pyta, czy zeszłomiesięczna zmiana modelu pomogła, czy
zaszkodziła. Dla inżyniera platformy, który już utrzymuje Prometheusa i chce
mieć w nim sygnały jakości. Dla programisty, który chce wiedzieć, która z jego
sesji poszła źle i w której turze. Dla organizacji, która musi obronić ten
pomiar przed radą pracowniczą, zanim w ogóle wolno mu zaistnieć.

Nie jest dla kogoś, kto chce tabeli wyników. To nie jest ograniczenie do
obejścia --- patrz rozdział~7.3.

\subsection{Jak wygląda sesja po zmierzeniu}

\begin{figure}[htbp]
\centering
\includegraphics[width=0.62\linewidth]{../diagrams/rendered/a1-system-context.pdf}
\caption{Granica systemu. Czterej asystenci emitują OpenTelemetry na jeden
endpoint; kolektor rozróżnia ich po zawartości, a nie po konfiguracji. Metrics
API GitHuba jest wyjątkiem --- jest odpytywane, a to, co zwraca, to zużycie na
poziomie organizacji: kontekst obok oceny, a nie sama ocena.}
\label{fig:context}
\end{figure}

% ============================================================================
\section{Architektura}
% ============================================================================

\subsection{Czym jest każdy projekt}

\begin{figure}[htbp]
\centering
\includegraphics[width=\linewidth]{../diagrams/rendered/a2-solution-layout.pdf}
\caption{Kolektor jest produktem. Wszystko inne jest widokiem na niego,
udogodnieniem deweloperskim albo opcją włączaną świadomie.}
\label{fig:layout}
\end{figure}

\begin{longtable}{@{}p{0.30\linewidth}p{0.44\linewidth}p{0.18\linewidth}@{}}
\toprule
\textbf{Projekt} & \textbf{Rola} & \textbf{Zależności} \\
\midrule
\endhead
\codepath{src/CopilotScope.Collector} & Przyjmowanie OTLP/HTTP z własnym
dekoderem protobuf, agregacja sesji, silnik jakości, analiza tur, potok
analizatorów, REST API, eksporter Prometheusa, trwałość & tylko Npgsql \\
\codepath{src/CopilotScope.Dashboard} & Interfejs Blazor Server: sesje, miernik
jakości, analiza tur, transkrypt, widoki zespołowe & brak \\
\codepath{src/CopilotScope.ServiceDefaults} & Wspólne jądro: autoinstrumentacja,
health, service discovery, odporność HTTP. Bez logiki domenowej &
OpenTelemetry, Microsoft.Extensions \\
\codepath{src/CopilotScope.AppHost} & Orkiestracja Aspire na potrzeby
developmentu: Postgres, pgAdmin, obie usługi & Aspire.Hosting \\
\codepath{src/CopilotScope.JudgeAgent} & Opcjonalny sędzia LLM: G-Eval, SPUR,
RAGAS, głęboka analiza tarcia, wykrywanie ukończenia zadania. Azure AI Foundry
albo dowolny endpoint zgodny z OpenAI & Azure.AI, Microsoft.Agents.AI \\
\codepath{src/CopilotScope.AgentForge} & Opcjonalne agenty person, oparte na
transkryptach objętych zgodą & Azure.AI, Microsoft.Agents.AI \\
\codepath{tools/CopilotScope.Seeder} & Fabrykowane sesje demo wysyłane do
działającego kolektora & brak \\
\codepath{tools/CopilotScope.LogImporter} & Transkrypty Claude Code
przetwarzane na ocenione sesje & brak \\
\codepath{tools/CopilotScope.TelemetryGen} & Jedna realistyczna sesja przez
prawdziwe OTLP/HTTP, używana jako sonda & brak \\
\bottomrule
\caption{Rozwiązanie projekt po projekcie. Trzy narzędzia jadą razem w jednym
obrazie kontenera \code{copilotscope-tools}, więc żadne z nich nie wymaga SDK
.NET do uruchomienia.}
\end{longtable}

Kolektor ma w czasie działania \emph{jedną} zależność NuGet, a dashboard nie ma
żadnej. To jest celowe: dekoder protobuf OTLP napisano w repozytorium, na
podstawie oficjalnych numerów pól z plików \code{.proto}, zamiast wciągać SDK
OpenTelemetry, a dashboard odpytuje REST API kolektora co dwie sekundy zamiast
nieść ze sobą framework JavaScript.

\subsection{Co dzieje się z jedną paczką telemetrii}

\begin{figure}[htbp]
\centering
\includegraphics[width=0.72\linewidth]{../diagrams/rendered/a3-ingest-pipeline.pdf}
\caption{Każda strzałka opuszczająca ścieżkę szczęśliwą loguje, dlaczego ją
opuściła. Ten diagram jest odpowiedzią na „stos działa, a dashboard jest
pusty''.}
\label{fig:ingest}
\end{figure}

Dwie decyzje w tym potoku warto wypowiedzieć wprost, bo obie są z gatunku tych,
które wyglądają na ważne dopiero po tym, jak gdzieś poszły źle.

\textbf{Redakcja następuje przed agregacją.} W trybie prywatności wartości
identyfikujące są usuwane z paczki, zanim cokolwiek dalej --- magazyn, zapis
migawki, eksporter Prometheusa --- w ogóle je zobaczy. Gwarancja przeżywa więc
zrzut bazy danych, zamiast zależeć od tego, czy każda ścieżka odczytu pamiętała,
żeby ją zastosować.

\textbf{Zdekodowana zawartość ma twardy limit.} Kestrel ogranicza rozmiar
skompresowanego żądania; gzip osiąga około 1000:1. Nieograniczona kopia
skompresowanego ciała to wektor wyczerpania pamięci, a na kolektorze bez
skonfigurowanego klucza jest osiągalna przed uwierzytelnieniem. Limit wynosi
64~MB na paczkę.

\subsection{Gdzie mieszkają dane i czego eksmisja nie oznacza}

Sesje żyją w dwóch miejscach naraz. Przyjęcie telemetrii oznacza sesję jako
brudną; proces w tle zapisuje zbiór brudnych raz na sekundę, więc nawał
telemetrii nie staje się nawałem zapisów. Pełny stan sesji trafia do Postgresa
jako migawka \code{jsonb} --- liczniki, próbki czasu do pierwszego tokenu,
statystyki narzędzi, agregaty per tura, zmiany, oceny, ogon zdarzeń i
przechwycony transkrypt --- z kilkoma kolumnami wyciągniętymi na wierzch do
odpytywania.

\begin{figure}[htbp]
\centering
\includegraphics[width=0.78\linewidth]{../diagrams/rendered/a4-session-lifecycle.pdf}
\caption{Magazyn w pamięci to ograniczony zbiór roboczy, a nie zakres Twojej
historii. Eksmisja nie jest usunięciem.}
\label{fig:lifecycle}
\end{figure}

Odczyty są obsługiwane z Postgresa, z nałożonymi na wierzch żywymi agregatami z
pamięci, więc API odpowiada na temat całego archiwum, podczas gdy sesja, do
której ktoś właśnie pisze, pozostaje aktualna. Bez łańcucha połączenia kolektor
degraduje się do samej pamięci i raportuje \code{durable: false}, zamiast
udawać coś innego --- i to jest też powód, dla którego \code{dotnet run} na
gołej maszynie działa zupełnie bez bazy.

% ============================================================================
\section{Jak wprowadzić dane}
% ============================================================================

\subsection{Instalacja od początku do końca}

\begin{figure}[htbp]
\centering
\includegraphics[width=0.70\linewidth]{../diagrams/rendered/b1-install-and-connect.pdf}
\caption{Dwie komendy plus jeden ręczny krok na asystenta, którego żaden skrypt
nie wykona za Ciebie.}
\label{fig:install}
\end{figure}

Na jednej maszynie nie ma czego deklarować: żadnego klucza do wygenerowania,
żadnej zmiennej środowiskowej do wyeksportowania, żadnego JSON-a do ręcznej
edycji. Każdy publikowany port wiąże się z \code{127.0.0.1}, Postgres nie
publikuje portu w ogóle i ufa swojemu nieopublikowanemu gniazdu, a klucz
przyjmowania jest pusty, co jest trybem otwartym kolektora. Poświadczenie na
pętli zwrotnej niczego nie chroni, a kosztuje krok konfiguracyjny w każdym
kliencie --- i to właśnie ten krok ludzie mylą.

\subsection{Dlaczego konfiguracja bywała najtrudniejszą częścią}

\begin{figure}[htbp]
\centering
\includegraphics[width=\linewidth]{../diagrams/rendered/b2-where-each-switch-lives.pdf}
\caption{Cztery powierzchnie, trzy rodzaje miejsc do ich skonfigurowania i
tylko jedno z tych miejsc przeżywa otwarcie nowego terminala.}
\label{fig:switches}
\end{figure}

Claude Code jest przypadkiem najlepszym. Blok \code{env} w pliku
\code{~/.claude/settings.json} działa w każdym terminalu, w każdym projekcie, a
także dla Claude Code uruchomionego wewnątrz VS Code. Znaczenie mają cztery
wartości, a dwie pierwsze to te, które ludzie pomijają:

\begin{itemize}[leftmargin=*]
    \item \codewrap{CLAUDE_CODE_ENABLE_TELEMETRY}=1 --- przełącznik główny. Bez
    niego nie jest eksportowane nic, niezależnie od reszty ustawień.
    \item \code{OTEL_LOGS_EXPORTER}=otlp --- to zdarzenia logów niosą sesję.
    Domyślna instalacja emituje metryki i zdarzenia, a nie emituje span-ów, więc
    konfiguracja ustawiająca tylko eksporter metryk daje sesję niemal pustą.
    \item \codewrap{OTEL_EXPORTER_OTLP_PROTOCOL}=http/protobuf oraz endpoint.
\end{itemize}

Przechwytywanie treści to tutaj trzy osobne zgody
(\codewrap{OTEL_LOG_USER_PROMPTS}, \codewrap{OTEL_LOG_ASSISTANT_RESPONSES},
\codewrap{OTEL_LOG_TOOL_DETAILS}), a standardowa zmienna OpenTelemetry GenAI,
która działa dla Copilot CLI, \emph{nie} jest czytana przez Claude Code. Czas
do pierwszego tokenu wymaga bety śledzenia, która jest dla tego asystenta
jedynym źródłem tego sygnału.

VS Code konfiguruje się przez ustawienia użytkownika i wymaga przeładowania
okna, bo rozszerzenie czyta konfigurację przy starcie. Copilot CLI czyta
zmienne środowiskowe i nic poza nimi. Cowork konfiguruje się we własnym
interfejsie aplikacji desktopowej, chce pełnej ścieżki \code{/v1/logs} zamiast
bazowego endpointu, wymaga planu Team lub Enterprise oraz uprawnień
administratora organizacji, i eksportuje wyłącznie zdarzenia logów.

\subsection{Ścieżka, która nie wymaga żadnej konfiguracji}

\begin{figure}[htbp]
\centering
\includegraphics[width=0.72\linewidth]{../diagrams/rendered/b3-transcript-import.pdf}
\caption{Claude Code zapisuje każdą sesję na dysk niezależnie od tego, czy
telemetria jest skonfigurowana --- a większość programistów nigdy jej nie
konfiguruje.}
\label{fig:import}
\end{figure}

Dwie odmowy na tej ścieżce są celowe i warte obrony. Etykieta repozytorium
zostaje \emph{pusta} zamiast zgadywana z nazwy katalogu, bo zgadnięcie
wymyśliłoby drugą kohortę dla projektu, który kolektor już zna po zdalnym
adresie gita --- dwa wiersze na jedno repozytorium, co jest gorsze niż brak
etykiety. A sesja, którą kolektor już ma z żywej telemetrii, nie jest
nadpisywana przez import, bo transkrypt niesie mniej sygnału i po cichu
obniżyłby ocenę zmierzoną na większym materiale.

Sesje zaimportowane są oznaczone i niosą realnie niższą pewność. Transkrypt
rejestruje tokeny, modele, wywołania narzędzi i ich wyniki, granice tur oraz
prawdziwe czasy zegarowe. Nie rejestruje czasu do pierwszego tokenu, decyzji o
przyjęciu lub odrzuceniu zmiany ani ocen kciukiem, bo to są zdarzenia
OpenTelemetry. Zamiast ustawiać je na zero --- co czytałoby się jako „każda
zmiana odrzucona'' --- są pozostawione jako brakujące, a wskaźnik złożony
renormalizuje.

% ============================================================================
\section{Pomiar}
% ============================================================================

\subsection{Wskaźnik złożony}

\begin{figure}[htbp]
\centering
\includegraphics[width=0.80\linewidth]{../diagrams/rendered/c1-composite-score.pdf}
\caption{Sześć składowych. Do wskaźnika wchodzą tylko te, które dostarczyły
dane, a wagi są renormalizowane po nich.}
\label{fig:composite}
\end{figure}

\begin{center}
\begin{tabular}{@{}llp{0.52\linewidth}@{}}
\toprule
\textbf{Składowa} & \textbf{Waga} & \textbf{Czym jest} \\
\midrule
Niezawodność & 0,25 & Kwadrat odsetka wywołań bez błędu, dla LLM i narzędzi.
Kwadrat sprawia, że sesja z kilkoma błędami jest karana bardziej niż liniowo \\
Akceptacja   & 0,20 & Przyjęte zmiany, sparowane z przeżywalnością zmian, żeby
przyjmowanie złych podpowiedzi nie czytało się jako sukces \\
Tarcie       & 0,20 & Średnia ocena tury z analizatora tur \\
Opóźnienie   & 0,15 & Jaka część sesji leżała powyżej progów 2~s uwagi i 8~s
porzucenia \\
Oceny        & 0,10 & Kciuki, tam gdzie asystent je raportuje \\
Efektywność  & 0,10 & Ekonomia tokenów: koszt na turę i na przyjętą zmianę \\
\bottomrule
\end{tabular}
\end{center}

\vspace{0.3cm}

\finding{Renormalizacja jest tym, co trzeba zrozumieć.} Do wskaźnika wchodzą
tylko składowe, które mają dane, a wagi są przeskalowane po nich. Wcześniejsza
wersja pozwalała pustym składowym wnosić neutralny prior, w efekcie czego
niemal każda sesja dostawała około 80 --- pomiar raportował prior, a nie sesję.
Pewność, eksportowana obok każdej oceny, to pokrycie danymi pomnożone przez
rampę liczby próbek: 90 zbudowane na czterech próbkach znaczy mniej niż 70
zbudowane na czterdziestu.

Oceny słowne są progami na wskaźniku: 85 i wyżej to doskonale, 70 dobrze, 55
przeciętnie, 40 słabo, poniżej --- krytycznie. Ranga percentylowa liczona jest
względem stałego okna wstecznego sesji użytkownika, a nie względem tego, co
akurat było w pamięci, i to właśnie czyni ją porównywalną w czasie.

\subsection{Analiza tur, czyli która tura poszła źle}

\begin{figure}[htbp]
\centering
\includegraphics[width=0.88\linewidth]{../diagrams/rendered/c2-turn-analysis.pdf}
\caption{Wskaźnik złożony mówi, że sesja była słaba. To mówi, która tura i
dlaczego.}
\label{fig:turns}
\end{figure}

Każdy trace \code{invoke_agent} to jedna tura. Każda tura jest oceniana pod
kątem błędów LLM i narzędzi, opóźnienia względem mediany czasu do pierwszego
tokenu \emph{tej właśnie sesji}, oraz pętli naprawczych --- serii wywołań
narzędzi zawierających niepowodzenia. Dashboard pokazuje najlepszą i najgorszą
turę wraz z powodami.

Mediana liczona per sesja jest tu wyborem nośnym. Globalny próg opóźnienia nie
odróżni wolnej tury od wolnego modelu ani wolnego modelu od wolnego poranka;
sesja porównana sama ze sobą --- owszem.

Warto zauważyć konsekwencję dla asystentów, które nie emitują granicy tury:
eksport Cursora wysyła metryki i logi, ale nie span-y, a tura \emph{jest} tu
trace'em, więc analiza tarcia na poziomie tur nie może dla niego zadziałać w
ogóle, nawet gdyby eksport działał.

\subsection{Potok analizatorów}

Analizatory są wtykowe: jedna klasa \code{IInsightAnalyzer} plus jedna
rejestracja w kontenerze wstrzykiwania zależności dodaje algorytm bez żadnej
pracy po stronie interfejsu. Te działające w procesie to przeżywalność zmian,
przepustowość, użyteczność opóźnienia, ekonomia tokenów i tarcie pracy.

Algorytmy oceniane przez LLM celowo \emph{nie} są w tym potoku. Interfejs jest
synchroniczny i wewnątrzprocesowy; wywołanie sędziego to asynchroniczne
wywołanie sieciowe do serwera modelu, z kosztem i limitem tempa. Mieszkają
zamiast tego w osobnej, opcjonalnej usłudze sędziego, dostępnej pod
\code{POST /api/sessions/{id}/judge}, a tę usługę można skierować na Azure AI
Foundry albo na endpoint zgodny z OpenAI na własnym sprzęcie --- więc wdrożenie
lokalne może je uruchomić bez wypuszczania transkryptów poza budynek.

% ============================================================================
\section{Dziesięć algorytmów i ile każdy z nich dziś znaczy}
% ============================================================================

Projekt śledzi dziesięć algorytmów oceny na dwóch osiach: czy implementacja
istnieje i czy działa lokalnie, czy wymaga dostępu do modelu.

\begin{longtable}{@{}p{0.03\linewidth}p{0.26\linewidth}p{0.20\linewidth}p{0.42\linewidth}@{}}
\toprule
\textbf{\#} & \textbf{Algorytm} & \textbf{Status} & \textbf{Uwagi} \\
\midrule
\endhead
1 & LLM jako sędzia (G-Eval) & zaimplementowany, opcjonalny & Wymaga endpointu
LLM; model lokalny wystarcza \\
2 & Rubryki satysfakcji SPUR & zaimplementowany, zero-shot & Kierunkowy, dopóki
nie powstaną oznaczone sesje do kalibracji \\
3 & Metryki składowych RAG (RAGAS) & zaimplementowany, opcjonalny & Sensowny
tylko dla przepływów z wyszukiwaniem; raportuje brak danych, dopóki kontekst
wyszukiwania nie jest przechwytywany \\
4 & Przeżywalność zmian & pełny, lokalny & Podział na cztero-gramy i brak
cofnięcia, 0,4/0,6; zasila też składową akceptacji \\
5 & Przepustowość ważona akceptacją & pełny, lokalny & Przyjęte linie na turę i
na 1k tokenów wyjścia, z dyskontem za odrzucenia \\
6 & Tarcie i naprawy na poziomie tur & pełny, lokalny & Analizator tur; także
składowa tarcia \\
7 & Model użyteczności opóźnienia & pełny, lokalny & Krzywa użyteczności per
próbka, z kubełkami ryzyka 2~s i 8~s \\
8 & Ekonomia tokenów i cache & pełny, lokalny & Koszt per model, oszczędność
cache, koszt na turę i na przyjętą zmianę \\
9 & Sygnały tarcia pracy & uproszczony lokalnie, głęboki w sędzim & Liczy
zaobserwowane zdarzenia naprawcze, a nie wnioskowane emocje. Domyślnie
wyłączony, tylko raportujący, nigdy nie oceniany \\
10 & Wykrywanie ukończenia zadania & częściowy lokalnie, zaimplementowany w
sędzim & Brak jeszcze lokalnej ścieżki wejścia dla kodów wyjścia builda i
testów \\
\bottomrule
\caption{Status implementacji. „Opcjonalny'' oznacza konfigurację, a nie to,
czy kontener działa.}
\end{longtable}

Algorytm 9 zasługuje na własne zdanie, bo jest tym, który najłatwiej źle
odczytać. Liczy \emph{zaobserwowane zdarzenia naprawcze} --- przeformułowania
mierzone podobieństwem Jaccarda między kolejnymi promptami, repliki korygujące,
markery negatywnej oceny --- a nie emocje. Jest celowo wykluczony ze wskaźnika
złożonego: heurystyka leksykalna nie należy do liczby, na podstawie której ktoś
mógłby działać. Jest też jedynym analizatorem czytającym własny tekst
programisty, dlatego pozostaje wyłączony, dopóki operator go nie włączy, i
dlatego cytowanie wiadomości wymaga drugiej, osobnej zgody.

% ============================================================================
\section{Pokrycie sygnałami: dlaczego 80 tutaj to nie 80 tam}
% ============================================================================

\begin{figure}[htbp]
\centering
\includegraphics[width=0.92\linewidth]{../diagrams/rendered/c3-signal-coverage.pdf}
\caption{Wszyscy asystenci lądują w tym samym schemacie. Nie raportują tych
samych sygnałów.}
\label{fig:coverage}
\end{figure}

Wszystkie cztery wspierane powierzchnie normalizują się do jednego
wewnętrznego słownika, więc agregacja, ocena i dashboard traktują je
identycznie. Różni się to, co każda z nich faktycznie wysyła.

\begin{center}
\begin{tabular}{@{}lcccc@{}}
\toprule
\textbf{Sygnał} & \textbf{VS Code} & \textbf{Copilot CLI} & \textbf{Claude Code}
& \textbf{Cowork} \\
\midrule
Tokeny, wywołania, błędy       & tak & tak & tak & tak \\
Czas do pierwszego tokenu      & tak & tak & tylko z betą & nie \\
Przyjęcie / odrzucenie zmiany  & tak & nie & tak & tak \\
Oceny kciukiem                 & tak & nie & nie & nie \\
Linie kodu                     & tak & nie & tak & nie \\
Granica tury (trace)           & tak & tak & tylko z betą & nie \\
\bottomrule
\end{tabular}
\end{center}

\vspace{0.3cm}

\finding{Oceny są porównywalne wewnątrz asystenta i kierunkowe pomiędzy nimi.}
Sesja Claude Code nie ma kciuków, a bez bety śledzenia nie ma czasu do
pierwszego tokenu, więc jej 80 opiera się na mniejszym materiale niż 80 sesji
VS Code. Wskaźnik złożony radzi sobie z tym uczciwie, renormalizując, a miara
pewności raportuje, na czym musiał pracować. Czego nie potrafi, to sprawić,
żeby obie liczby znaczyły to samo --- więc porównanie asystentów musi utrzymać
stały zestaw sygnałów, a nie tylko wskaźnik.

Cursor jest przypadkiem szczególnym i jako taki jest udokumentowany. Jego
eksport OpenTelemetry dostępny jest wyłącznie w planie Enterprise i wysyła
metryki oraz logi, ale nie span-y. Kolektor wykrywa i normalizuje atrybuty
\code{cursor.*}, ale żadna zawartość z prawdziwej sesji Cursora nigdy nie
została przeciw niemu przetestowana, więc jest wymieniony jako wykrywany i
niezweryfikowany, a nie jako wspierany. To rozróżnienie istnieje, bo ten projekt
już raz zadeklarował wsparcie dla asystenta, którego zawartości nigdy nie
widział.

% ============================================================================
\section{Uczciwe ograniczenia}
% ============================================================================

Ten rozdział istnieje, bo pomiar podany bez swoich ograniczeń jest gorszy niż
brak pomiaru: czytelnik mu ufa.

\subsection{Wskaźnik złożony nie jest skalibrowany}

W tym repozytorium nie ma żadnych etykiet ludzkich i nie przeprowadzono żadnej
kalibracji. Wagi są przemyślanym osądem, a nie dopasowanym modelem. Każda ocena
sędziego jest więc kierunkowa i niczego nie blokuje.

Maszyneria, która to naprawi, istnieje i czeka na wejście. Usługa sędziego
niesie moduł kalibracji liczący $\kappa$ Cohena --- nieważone oraz ważone
liniowo i kwadratowo, każde z zaszczepionym przedziałem ufności bootstrap ---
pomiędzy sędzią a panelem ludzi, plus zgodność panelu z samym sobą jako sufit,
ponad który żaden algorytm nie wejdzie. Werdykty to \emph{skalibrowany},
\emph{nieskalibrowany}, \emph{sufit za nisko} albo \emph{za mało danych},
względem konfigurowalnego progu, domyślnie $\kappa \geq 0{,}70$.

Brakowało dotąd czegokolwiek do skonsumowania. Przepływ etykietowania istnieje
teraz za flagą \codewrap{CopilotScope:Labelling:Enabled}: panel oceniającego na
widoku sesji, z uchwytem oceniającego, pasmem od 0 do 3 na rubrykę wraz z
tekstem kotwiczącym i opcją pominięcia, która jest zapisywana jako prawdziwa
odpowiedź, ale nie jest eksportowana jako etykieta.
\code{GET /api/labels/export} produkuje dokładnie ten JSON, który czyta silnik
kalibracji, przypięty testem obiegu w obie strony.

Sesje zasiane są domyślnie wykluczone z eksportu i to wykluczenie jest istotą
rzeczy: walidowanie modelu oceny względem syntetycznych person napisanych przez
to samo repozytorium byłoby kołem, a $\kappa$ zbudowane na kole jest gorsze niż
brak $\kappa$.

\subsection{Pojedyncza liczba nie jest werdyktem}

Pewność jest eksportowana obok każdej oceny nie bez powodu. Czytaj składowe, a
nie nagłówek. Sesja z wynikiem 62, wysoką pewnością i jasno wskazaną najgorszą
turą jest bardziej użyteczna niż 88 na dwóch próbkach.

\subsection{To nie jest tabela wyników programistów --- i to jest egzekwowane}

Ocena dotyczy \emph{sesji}, a nie osoby. Rankingowanie programistów według niej
jest nadużyciem, którego projekt nie wspiera: w filtrze kohort nie ma wymiaru
per programista, nie ma takiego widoku ani eksportu z taką kolumną. Testy tego
pilnują.

We wdrożeniu zespołowym konwencja staje się maszynerią. Tryb prywatności
pseudonimizuje tożsamości przy przyjęciu, odrzuca treść promptów i odpowiedzi,
odmawia wyrenderowania widoku obejmującego mniej niż $k$ programistów i loguje
każdy odczyt.

Uzasadnienie nie jest wyłącznie etyczne. Ocena używana do oceny pracowniczej
niemal natychmiast przestaje mierzyć narzędzie, bo mierzeni ludzie się
adaptują --- a adaptują to, co jest liczone, a nie to, co chciałeś zmierzyć. To
jest też dokładnie powód, dla którego wskaźnik akceptacji to tylko 0,20
wskaźnika złożonego i jest sparowany z przeżywalnością zmian: naciskaj na samą
akceptację, a nagrodzisz przyjmowanie złych podpowiedzi. Prawo Goodharta stosuje
się do tego repozytorium tak samo jak do wszystkiego, co ono mierzy.

\subsection{Czego projekt celowo nie robi}

Nie archiwizuje surowej telemetrii --- przechwytywanie treści zapisuje
ograniczony podgląd, ostatnie sto wpisów przyciętych do czterech tysięcy znaków,
a pełne archiwum to zadanie dla prawdziwego backendu telemetrycznego, który
karmi forwarder. Nie wlicza dostarczonych rezultatów do oceny, mimo że wiąże już
sesje z pull requestami, które wyprodukowały, bo badanie korelacji, które by to
uzasadniło, jest właśnie powodem zbierania tych rezultatów. I nie promuje
analizatora tarcia do wskaźnika złożonego, z powodu podanego w rozdziale~5.

% ============================================================================
\section{Prywatność, bezpieczeństwo i wdrożenia zespołowe}
% ============================================================================

\subsection{Moment, w którym pojawia się telemetria drugiego programisty}

W chwili, gdy sesje drugiej osoby trafiają do tego samego kolektora, prowadzisz
system techniczny zdolny do monitorowania wydajności pracowników. W Unii
Europejskiej uruchamia to współdecydowanie rady pracowniczej na podstawie samej
zdolności, niezależnie od Twoich zamiarów. Tryb prywatności jest tym, co czyni
intencję egzekwowalną, a nie tylko zadeklarowaną.

\begin{itemize}[leftmargin=*]
    \item Tożsamości są pseudonimizowane, zanim cokolwiek je zapisze, solą,
    którą dostarczasz. Sól ulotna jest generowana, jeśli o niej zapomnisz, a
    kolektor ostrzega głośno, bo pseudonimy zmieniające się przy każdym
    restarcie przestają korelować historię między wdrożeniami.
    \item Treść promptów i odpowiedzi jest odrzucana przy przyjęciu, niezależnie
    od tego, jak skonfigurowani są klienci.
    \item Żaden widok nie renderuje się dla mniej niż $k$ programistów.
    \item Każdy odczyt jest logowany do tabeli, której nie dotyka zamiatanie
    retencji sesji, bo zapis audytowy musi przeżyć sesje, które opisuje.
    \item Przekazywanie surowego OTLP jest odrzucane, chyba że operator wprost
    stwierdzi, że backend docelowy jest objęty tą samą umową --- wierne
    przekazanie niezredagowanej zawartości to dziura przez wszystkie pozostałe
    zabezpieczenia.
\end{itemize}

\code{GET /api/privacy} raportuje, co działające wdrożenie faktycznie
egzekwuje, więc odpowiedź na pytanie „czy jest włączony?'' nie sprowadza się do
zaufania plikowi konfiguracyjnemu.

\subsection{Model poświadczeń}

\begin{figure}[htbp]
\centering
\includegraphics[width=0.80\linewidth]{../diagrams/rendered/d1-deployment-modes.pdf}
\caption{Domyślny brak poświadczeń jest ograniczony zakresem, a nie
rozluźniony. Publikacja poza pętlę zwrotną to inne wdrożenie i tak jest
traktowana.}
\label{fig:deploy}
\end{figure}

Bramkowanie idzie za obecnością klucza, w każdym środowisku: przy pustym
\codewrap{CopilotScope:Ingest:ApiKey} kolektor działa otwarcie, a przy kluczu
ustawionym przyjmowanie, API zapytań i endpoint Prometheusa wymagają go
wszystkie.

Dla czegokolwiek współdzielonego podziel ten jeden klucz na zakresy. Klucz,
który trzyma edytor każdego programisty, nie powinien być zarazem kluczem
czytającym transkrypty i kasującym historię:

\begin{itemize}[leftmargin=*]
    \item \textbf{Ingest} --- \code{POST /v1/*} i nic poza tym. Trzymany przez
    każdy emiter.
    \item \textbf{Read} --- API zapytań i \code{/metrics}, co obejmuje
    przechwycone transkrypty. Trzymany przez dashboard i Prometheusa.
    \item \textbf{Admin} --- operacje destrukcyjne i administracyjne. Zawiera w
    sobie Read.
\end{itemize}

Dashboard ma własne logowanie z dwiema rolami, bo transkrypty są wrażliwym
ładunkiem: widz widzi oceny, tury i agregaty; administrator widzi także treść
rozmów i może usuwać. Żaden plik compose nie terminuje TLS, więc cokolwiek
współdzielonego potrzebuje reverse proxy przed sobą --- te poświadczenia
podróżują w nagłówku i w ciasteczku.

Kolektor nie jest w stanie sam ustalić, czy jest wystawiony: za Dockerem każde
żądanie przychodzi z bramy mostka, więc adres zdalny niczego nie dowodzi. Pliki
compose przekazują mu zatem adres, pod którym go opublikowano, a ostrzeżenie
startowe pojawia się, gdy to nie jest pętla zwrotna, a klucza nie ma. Dwie
komendy odmawiają takiego połączenia wprost, zamiast o nim ostrzegać.

% ============================================================================
\section{Samouczek}
% ============================================================================

Wszystko poniżej da się uruchomić. Jedynym wymaganiem jest Docker; SDK .NET jest
potrzebne wyłącznie na ścieżce dla kontrybutora w punkcie~9.9.

\subsection{Pięć minut: od pustej maszyny do ocenionej sesji}

\subsubsection*{Krok 1 --- instalacja}

\begin{shell}
curl -fsSL https://raw.githubusercontent.com/konradcinkusz/\
copilot-scope/master/install.sh | sh
\end{shell}

Na Windowsie, w PowerShellu:

\begin{shell}
irm https://raw.githubusercontent.com/konradcinkusz/copilot-scope/\
master/install.ps1 | iex
\end{shell}

Instalator sprawdza Dockera, zapisuje plik compose oraz skrypt sterujący
\code{copilotscope} do katalogu \code{~/.copilotscope}, uruchamia stos, czeka aż
kolektor odpowie i proponuje skonfigurowanie każdego znalezionego asystenta.
Dodaj \code{--yes} po \code{| sh -s --}, żeby przyjąć wszystkie propozycje bez
pytań.

Nic więcej nie jest wymagane. Nie ma klucza do wygenerowania ani pliku
\code{.env} do wypełnienia.

\subsubsection*{Krok 2 --- udowodnij, że potok działa, zanim obwinisz klienta}

\begin{shell}
copilotscope probe
\end{shell}

To wysyła jedną symulowaną sesję prawdziwym OTLP/HTTP protobuf, a potem
sprawdza, czy kolektor potrafi ją oddać. Jeśli przejdzie, działają
przyjmowanie, dekodowanie, ocena i trwałość, a wszystko, czego nadal brakuje,
jest konfiguracją klienta. Zrobienie tego najpierw oszczędza godzinę
debugowania niewłaściwej połowy systemu.

\subsubsection*{Krok 3 --- zobacz cokolwiek}

\begin{shell}
copilotscope demo          # sesje fabrykowane, oznaczone jako demo
copilotscope open          # http://localhost:5200
\end{shell}

Sesje zasiane są wyraźnie oznaczone, więc zrzut ekranu z nich nigdy nie zostanie
wzięty za dowód na temat prawdziwego asystenta. \code{copilotscope demo demo}
ładuje większy, wielodniowy zestaw zamiast szybkiego.

\subsection{Podłączanie każdego asystenta}

\begin{shell}
copilotscope connect claude-code     # ~/.claude/settings.json
copilotscope connect vscode          # ustawienia użytkownika VS Code
copilotscope connect copilot-cli     # rc powłoki lub zakres User na Windows
copilotscope connect cowork          # wypisuje, co wpisać w aplikacji
copilotscope connect all             # Claude Code i VS Code razem
\end{shell}

Dwie flagi działają dla wszystkich:

\begin{itemize}[leftmargin=*]
    \item \code{--capture} eksportuje także tekst promptów i odpowiedzi.
    Domyślnie wyłączone; to przełącznik, który zamienia wdrożenie trzymające
    same metadane we wdrożenie trzymające kod źródłowy i wszystko, co
    programista wkleił do czatu.
    \item \code{--print} pokazuje dokładnie, co zostałoby zapisane, i niczego
    nie zmienia.
\end{itemize}

\code{--traces} dotyczy tylko Claude Code i włącza betę śledzenia, jedyne
źródło czasu do pierwszego tokenu dla tego asystenta. Schemat span-ów może się
jeszcze zmienić; to właśnie znaczy flaga bety.

Potem ręczny krok, którego narzędzia nie wykonają za Ciebie: przeładuj okno VS
Code, zrestartuj Claude Desktop dla Cowork albo otwórz nowy terminal dla
Copilot CLI. Na koniec rozmawiaj w \emph{trybie agenta} --- same podpowiedzi
inline nie produkują telemetrii czatu, co jest drugim najczęstszym powodem
pustego dashboardu.

\code{copilotscope disconnect <cel>} usuwa dokładnie te klucze, które zostały
dodane, zostawiając resztę ustawień nietkniętą.

\subsection{Ocena historii, którą już masz}

Jeśli używasz Claude Code, pomiń wszystko powyżej i uruchom:

\begin{shell}
copilotscope import --dry-run     # zobacz co znalazł, nie wysyłaj nic
copilotscope import               # zaimportuj
\end{shell}

Ponowne uruchomienie jest bezpieczne: sesje zachowują własny identyfikator
Claude Code, więc drugi przebieg zastępuje, a nie duplikuje. Można to wstawić do
harmonogramu. Tekst promptów zostaje poza importem, chyba że podasz
\code{--include-content}.

Jedno zastrzeżenie, które sama komenda wypisuje: uruchomiony w kontenerze
importer nie sięgnie do gita, więc zaimportowane sesje nie niosą etykiety
repozytorium. Kiedy to ma znaczenie, uruchom go z klona z SDK .NET.

\subsection{Jak czytać sesję}

Otwórz sesję i czytaj ją w tej kolejności.

\begin{enumerate}[leftmargin=*]
    \item \textbf{Ocena i pewność razem.} Nigdy sama ocena. Niska pewność
    oznacza, że sesja nie wyemitowała dość, by liczba miała wagę.
    \item \textbf{Najgorsza tura wraz z powodami.} To część, na której można
    działać: błąd, wartość odstająca opóźnienia względem własnej mediany tej
    sesji, albo pętla naprawcza.
    \item \textbf{Składowe.} Wskaźnik 62 napędzany opóźnieniem to inny problem
    niż 62 napędzane niezawodnością, i mają różne naprawy.
    \item \textbf{Przeżywalność zmian, jeśli asystent ją raportuje.} Akceptacja
    bez przeżywalności to sygnatura kodu przyjętego i zaraz cofniętego.
    \item \textbf{Ekonomia tokenów.} Koszt na przyjętą zmianę to liczba, która
    zamienia rozmowę o jakości w rozmowę o budżecie.
\end{enumerate}

Wbudowana strona dokumentacji pod \code{/docs} objaśnia każdy kafelek i każdą
składową językiem interpretacji, a nie językiem metryk.

\subsection{Porównywanie okien i kohort}

Dwa pytania mają własne endpointy, oba filtrowalne po repozytorium, asystencie i
modelu, oba eksportowalne do CSV wyłącznie z wierszami grup:

\begin{shell}
GET /api/cohorts?days=30      # zwiń okno po każdej osi
GET /api/compare?days=30      # to okno względem poprzedniego
GET /api/digest?days=7        # tydzień, jako coś do przesłania dalej
\end{shell}

\code{/api/compare} to pytanie „czy zmiana pomogła'', z bazą odniesienia
domyślnie ustawioną na równie długie okno bezpośrednio poprzedzające bieżące.

Każda oś opisuje narzędzie. Żadna nie opisuje programisty i nie jest to
przeoczenie.

\subsection{Prometheus i Grafana}

Jeśli już je utrzymujesz, kolektor wystawia swoje sygnały \emph{wyliczone} ---
nie samo surowe zużycie --- w formacie tekstowym Prometheusa:

\begin{shell}
docker compose -f docker-compose.grafana.yml up -d
# Grafana na http://localhost:3000, datasource i dashboard zaprowizjonowane
\end{shell}

Oceny są eksportowane jako pary \code{_sum}/\code{_count}, a nie jako
uśrednione wcześniej gauge'e, więc PromQL utrzymuje uczciwą arytmetykę przy
dowolnym grupowaniu:

\begin{shell}
sum by (emitter) (copilotscope_quality_score_sum)
  / sum by (emitter) (copilotscope_quality_score_count)
\end{shell}

Wszystko niesie etykietę \code{emitter}, więc regresja w jednym asystencie jest
widoczna, a nie uśredniona na nic. Serie per sesja są opcjonalne i ograniczone
limitem, bo identyfikatory sesji są nieograniczone i wręczyłyby Prometheusowi
nową serię czasową dla każdej odbytej rozmowy.

Jedna pułapka zasługuje na nazwanie. Zaprowizjonowane reguły alertów używają
ilorazów gauge'ów i nigdy \code{rate()} na liczniku. Kolektor przelicza swoje
rodziny Prometheusa z ograniczonego zbioru w pamięci, więc rodzina
zadeklarowana jako licznik może zmaleć, gdy sesje wypadają ze zbioru, a
\code{rate()} czyta spadek jako reset licznika --- reguła zbudowana w ten
sposób odpala na eksmisji, a nie na czymkolwiek rzeczywistym.

\subsection{Alerty i tygodniowy skrót}

Dashboard, do którego trzeba wejść, zostaje porzucony. Wyjście, które wyzwala
decyzję, zostaje odnowione. Skonfiguruj \code{CopilotScope:Alerts}, a kolektor
porówna każdy tydzień z poprzednim, po repozytorium, asystencie i modelu, i
wyśle webhooka, gdy średni wskaźnik kohorty spadnie.

To, co celowo \emph{nie} odpala, jest tu częścią ciekawszą: spadek, któremu
towarzyszył spadek \emph{pewności}, jest raportowany jako zmiana podstawy
pomiaru, a nie jako regresja. Wskaźnik renormalizuje po składowych, które mają
dane, więc kohorta, która przestała raportować sygnał, jest mierzona inaczej, a
nie gorzej. Wysyłanie zespołu w pogoń za zmianą, która nigdy nie zaszła, to
sposób, w jaki wycisza się kanał alertów, a wyciszony kanał jest wart mniej niż
brak kanału, bo zespół wierzy, że ma pokrycie.

\subsection{Zamiana oceny z opinii w pomiar}

To najwartościowsza rzecz, jaką czytelnik tego podręcznika może wnieść.

\begin{shell}
# appsettings.json kolektora
"CopilotScope": { "Labelling": { "Enabled": true } }
\end{shell}

Na widoku sesji pojawia się panel \emph{Oceń tę sesję}. Oceniaj sesje
prawdziwe, nie zasiane. Następnie:

\begin{shell}
GET /api/labels/export > calibration/labels.json
POST /api/calibration/report   # na agencie sędziego, :5400
\end{shell}

Endpoint raportu to czysta arytmetyka: deterministyczna, darmowa i uruchamialna
w CI. Plan badawczy prosi o pięćdziesiąt do stu oznaczonych sesji, zanim
$\kappa$ zacznie cokolwiek znaczyć.

\subsection{Praca z klona}

\begin{shell}
git clone https://github.com/konradcinkusz/copilot-scope.git
cd copilot-scope
dotnet build
dotnet run --project src/CopilotScope.AppHost     # Aspire: wszystko
dotnet test                                        # Docker niepotrzebny
\end{shell}

Dwie opcjonalne usługi agentowe siedzą za profilem compose, więc zwykłe
\code{docker compose up} ich nie uruchamia:

\begin{shell}
docker compose --profile agents up -d --build
\end{shell}

\subsection{Kiedy nic się nie pojawia}

\begin{shell}
copilotscope doctor
\end{shell}

Sprawdza Dockera i kontenery, dokument zdrowia kolektora, czy wdrożenie jest
wystawione bez klucza, co faktycznie mówi plik ustawień każdego asystenta, czy
wyeksportowana zmienna \codewrap{OTEL_EXPORTER_OTLP_ENDPOINT} nie nadpisuje tego
pliku, oraz ile transkryptów leży na dysku niezaimportowanych.

Klasyczne przyczyny, w kolejności rzeczywistego występowania:

\begin{enumerate}[leftmargin=*]
    \item VS Code nie został przeładowany po zmianie ustawień.
    \item Same podpowiedzi inline --- nie było tury czatu ani agenta.
    \item Dla Claude Code: brakuje przełącznika głównego albo eksportera logów.
    \item Wyeksportowana zmienna środowiskowa wygrywa z plikiem ustawień dla
    wszystkiego, co uruchomiono z tej powłoki.
    \item Korporacyjne ustawienia zarządzane przypinają endpoint do firmowego
    kolektora. Wartości zarządzane zawsze wygrywają.
    \item Cowork wskazano na bazowy endpoint zamiast na pełną ścieżkę
    \code{/v1/logs}, albo aplikacji nie zrestartowano.
\end{enumerate}

% ============================================================================
\section{Eksploatacja}
% ============================================================================

\subsection{Co blokuje merge, a co wydaje tag}

\begin{figure}[htbp]
\centering
\includegraphics[width=\linewidth]{../diagrams/rendered/d2-ci-gates.pdf}
\caption{Przez większość życia tego projektu obrazy budowano wyłącznie na tagu
wydania. To jest luka, którą zamyka zadanie budujące kontenery.}
\label{fig:ci}
\end{figure}

Zadanie budujące kontenery zapracowało na swoje miejsce już w pierwszym
przebiegu: znalazło dwa istniejące wcześniej błędy, przez które żaden obraz nie
dawał się zbudować. Pierwszym był wspólny projekt nieobecny w kontekście builda,
czego build całego rozwiązania nie mógł zobaczyć, bo tam projekt leży dokładnie
tam, gdzie wskazuje referencja. Drugi, osiągalny dopiero po naprawie
pierwszego, to dwa projekty webowe rozwiązujące swój \code{appsettings.json} do
tej samej ścieżki publikacji. Żaden nie jest egzotyczny; oba były niewidoczne,
bo nic nie budowało obrazów, dopóki wydanie nie było już wycinane.

\subsection{Retencja i to, co rośnie}

Historia rośnie, dopóki nie ustawisz polityki. Retencja jest domyślnie
wyłączona:

\begin{shell}
"CopilotScope": { "History": {
  "RetentionDays": 0,        // 0 = trzymaj wszystko; N = usuwaj bezczynne N+ dni
  "BaselineDays": 30,        // okno, po którym liczona jest ranga percentylowa
  "BaselineMaxSamples": 5000
} }
\end{shell}

\code{BaselineDays} jest tym, co czyni percentyl porównywalnym w czasie: ranga
liczona jest względem stałego okna wstecznego sesji użytkownika, a nie
względem tego, co akurat zostało w pamięci w chwili spojrzenia.

\subsection{Archiwizacja liczb samego dostawcy, zanim wygasną}

Metrics API Copilota od GitHuba wystawia okno 28 dni i nic starszego. Przy
skonfigurowanym \codewrap{CopilotScope:VendorMetrics} --- organizacja lub
enterprise, token tylko do odczytu i Postgres --- kolektor odpytuje codziennie i
zatrzymuje okno bezterminowo, kluczując per dzień, więc ponowne odpytanie
nadpisuje, zamiast gromadzić duplikaty.

Pełny dokument odpowiedzi jest zapisywany dosłownie obok wyciągniętych
liczników, bo dostawca kasuje oryginał: parser trzymający tylko to, co rozumiał
w danym dniu, po cichu wyrzuciłby historię, której nie da się pobrać ponownie.

Token potrzebuje dokładnie jednego zakresu odczytu i niczego więcej. Nic w tej
funkcji nie zapisuje, a token z zakresami zapisu ma większy promień rażenia niż
funkcja ma zastosowań. Wyłącznie poziom organizacji i zespołu: API można
poprosić o rozbicie per programista, a to celowo o nie nie prosi.

To jest kontekst, a nie pomiar. Liczenie zużycia nadal nie mówi, czy narzędzie
pomaga. Archiwum kupuje mianownik --- zdanie „liczba stanowisk wzrosła o
czterdzieści procent w marcu'' jest tym, co czyni trend jakości czytelnym ---
oraz coś użytecznego pierwszego dnia dla organizacji, która ma stanowiska i nie
ma nigdzie instrumentacji.

\subsection{Czy kod faktycznie pojechał?}

Każdy sygnał powyżej kończy się na granicy sesji. Sesja może przebiec czysto,
mieć przyjęte zmiany i mimo to wyprodukować zmianę, której nikt nie zmergował
--- co jest tą samą krytyką, którą ten projekt kieruje pod adresem liczników
zużycia, obróconą przeciw niemu samemu.

Wiązanie rezultatów zamyka tę pętlę, łącząc sesje z pull requestami. Skieruj
webhooka GitHuba na \code{POST /api/outcomes/github} z sekretem; dostarczenia
są weryfikowane HMAC-iem po surowym ciele, a bez sekretu endpoint nie jest w
ogóle zamapowany, bo zapisuje do tych samych danych, względem których ocena ma
być zaraz walidowana.

Złączenie jest heurystyką --- telemetria niesie repozytorium i gałąź, nigdy
commit --- więc każde powiązanie ma swoją pewność i powód, a dopasowanie po
samym repozytorium jest pokazywane, ale wykluczone z jakiejkolwiek korelacji.
Rezultaty stoją \emph{obok} oceny i celowo nie są do niej wliczane, z powodu
podanego w punkcie~7.4.

% ============================================================================
\section{Gdzie co leży}
% ============================================================================

\begin{longtable}{@{}p{0.42\linewidth}p{0.54\linewidth}@{}}
\toprule
\textbf{Ścieżka} & \textbf{Co tam jest} \\
\midrule
\endhead
\codepath{install.sh}, \codepath{install.ps1} & Instalator jednokomendowy \\
\codepath{scripts/copilotscope} & Skrypt sterujący: up, connect, import, demo,
probe, doctor \\
\codepath{docker-compose.ghcr.yml} & Kanoniczny stos „pobierz i uruchom'' \\
\codepath{docs/TUTORIAL.md} & Ręczne przejście krok po kroku, per asystent, plus
rozwiązywanie problemów \\
\codepath{docs/SIGNAL_COVERAGE.md} & Pełna macierz sygnałów per emiter \\
\codepath{docs/PRIVACY.md} & Mapa danych z art. 30 RODO, retencja, wzór aneksu
do porozumienia zakładowego \\
\codepath{docs/CALIBRATION.md} & Metoda, kotwice skali, jak czytać zły wynik \\
\codepath{docs/DIAGRAMS.md} & Każdy diagram z tego podręcznika, renderowany
przez GitHuba \\
\codepath{docs/diagrams/} & Te same diagramy jako źródła do ponownego użycia \\
\codepath{docs/papers/} & Ten podręcznik, oba wydania językowe \\
\codepath{docs/tutorials/} & Samouczki krok po kroku, w obu językach \\
\codepath{docs/architecture/} & Zapisy decyzji architektonicznych i przeglądy \\
\codepath{SECURITY.md} & Model zaufania, zakres po zakresie \\
\codepath{GOVERNANCE.md} & Kto to utrzymuje, co jest stabilne, co się stanie,
jeśli opiekun przestanie \\
\codepath{research/articles/} & Formalne opracowanie modelu oceny i dziesięć
tematów prac dyplomowych z niego wyprowadzonych \\
\bottomrule
\caption{Mapa. Tam, gdzie ten podręcznik i inny dokument się nie zgadzają, ten
drugi jest źródłem prawdy, a ten jest nieaktualny.}
\end{longtable}

% ============================================================================
\section{Podsumowanie}
% ============================================================================

Teza tego projektu jest wąska i warto ją powtórzyć precyzyjnie. Liczenie, ile
asystenta AI zużyto, nie mówi, czy pomógł. Ocena sesji --- owszem: niedoskonale,
na materiale różniącym się między asystentami, wzorem jeszcze niezwalidowanym
względem ludzkiego osądu i nigdy o osobie.

Trzy produkty komercyjne oceniają dziś sesje z asystentami programistycznymi, co
jest najmocniejszym zewnętrznym dowodem na słuszność tezy: trzy firmy z
nieporównanie lepszym rozpoznaniem rynku uznały, że problem jest prawdziwy. To,
co pozostaje tu wyróżniające, to nie pomysł, lecz warunki. Nic nie opuszcza
Twojej infrastruktury. Wzór jest jawny i wyprowadzalny ręcznie. Narzędzie nie
potrafi rankingować programistów, a to jest egzekwowane testami, a nie obiecane
w dokumencie.

Uczciwym następnym krokiem jest kalibracja, a ta potrzebuje oznaczonych sesji, a
nie kolejnego kodu. Maszyneria jest zbudowana, format eksportu przypięty testem,
a panel oddalony o jedną flagę konfiguracyjną. Do tego czasu wskaźnik złożony
jest opinią z przedziałem ufności --- co jest rzeczą użyteczną, pod warunkiem że
nikt nie czyta jej jako pomiaru.

\section*{Odtwarzalność}

\begin{shell}
git clone https://github.com/konradcinkusz/copilot-scope.git
cd copilot-scope
npm ci                    # przypięte CLI mermaida
npm run render:diagrams   # docs/diagrams/*.mmd -> wektorowe PDF-y
pdflatex docs/papers/copilotscope-manual.pl.tex   # dwa razy, dla spisu treści
\end{shell}

Oba wydania językowe buduje razem
\code{.github/workflows/build-docs-pdf.yml}, na żądanie. PDF-y są wynikiem
builda i nigdy nie są commitowane. \code{scripts/check-diagrams.mjs} przerywa
build, jeśli diagram rozejdzie się ze swoją osadzoną kopią, a
\code{scripts/check-doc-parity.mjs} przerywa go, jeśli jedno wydanie językowe
zostanie zmienione bez drugiego.

\end{document}
